\subsection{逼近定理}

在本小节中，A 是一个复交换幺 Banach 代数。

命题 4. --- 设 L 是 \(\mathbf{C}^{n}\) 的一个多项式凸紧子集，U 是 L 的一个开邻域。对每个函数 \(f\in\mathcal{O}(U;A)\)，存在一个由 \(\mathbf{C}^{n}\) 上以 A 为系数的多项式函数组成的序列，它收敛于 \(f|L\)（依 \(\mathcal{C}(L;A)\) 中的范数）。

可以假设 L 非空且 A 非零。设 P（相应地，\(P_{0}\)）是由 \(C^{n}\) 上以 A 为系数（相应地，以 C 为系数）的多项式函数在 L 上的限制所组成的集合。设 B（相应地，\(B_{0}\)）是 P（相应地，\(P_{0}\)）在

\(\mathcal{C}(\mathrm{L};\mathrm{A})\) 中的 Banach 代数闭包。用 \(\iota\) 表示 A 到 B 中由常值函数所组成的赋范子代数上的单射。

设 \(z_{1},\ldots,z_{n}\) 是 \(\mathbf{C}^{n}\) 的坐标函数在 L 上的限制；它们都是 \(\mathrm{B}_{0}\) 的元素，并且令 \(\boldsymbol{z}=(z_{1},\ldots,z_{n})\)，则依 I, p.~47, 命题 15 有 \(\mathrm{Sp}_{\mathrm{B}_{0}}^{n}(\boldsymbol{z})=\mathrm{L}\)。

设 \(f \in \mathcal{O}(\mathrm{U};\mathrm{A})\)。通过与 \(\iota\) 复合，函数 \(f\) 定义了一个元素 \(f_{\mathrm{B}} = \iota \circ f\)（属于 \(\mathcal{O}(\mathrm{U};\mathrm{B})\)）。由于 \(\mathrm{Sp}_{\mathrm{B}}^{n}(\boldsymbol{z}) \subset \mathrm{Sp}_{\mathrm{B}_{0}}^{n}(\boldsymbol{z}) \subset \mathrm{U}\)，可以构造 B 的元素 \(b = \Theta_{\boldsymbol{z}}(f_{\mathrm{B}})\)。设 \(\boldsymbol{w} = (w_{1},\ldots ,w_{n}) \in \mathrm{L}\)，并设 \(\varphi\) 是从 B 到 A 中的连续保单位元态射 \(g \mapsto g(\boldsymbol{w})\)。我们有 \(\varphi \circ \iota = \mathrm{Id}_{\mathrm{A}}\)，因而 \(\varphi \circ f_{\mathrm{B}} = f\)。由于 \(\varphi (z_i) = w_i\)，I, p.~66, 命题 3 蕴涵 \(\varphi (\Theta_z(f_{\mathrm{B}})) = \Theta_w(\varphi \circ f_{\mathrm{B}})\)。因此我们有

\[
b (\pmb {w}) = \varphi (b) = \varphi (\Theta_ {z} (f _ {\mathrm{B}})) = \Theta_ {\pmb {w}} (\varphi \circ f _ {\mathrm{B}}) = \Theta_ {\pmb {w}} (f) = f (\pmb {w}),
\]

（由 I, p.~65, 命题 2 的推论）。于是，我们有 \(f|_{\mathrm{L}} = b\)；特别地，\(f|_{\mathrm{L}}\) 属于 B。这就证明了该命题。

定理 2 (Oka--Weil). --- 设 K 是 \(C^{n}\) 的一个多项式凸紧子集，P 是由 \(C^{n}\) 上以 A 为系数的多项式函数在 K 的某个邻域中的芽所组成的集合。则 P 在 \(\mathcal{O}(K;A)\) 中稠密。更确切地说，\(\mathcal{O}(K;A)\) 的每个元素都是 P 中元素所成序列的极限。

考察 \(\mathcal{O}(\mathrm{K};\mathrm{A})\) 的一个元素，即某个函数 \(f\in \mathcal{O}(\mathrm{U};\mathrm{A})\) 的芽，其中 U 是 K 的一个开邻域。由 I, p.~48, 引理 7，存在 K 的一个包含于 U 且为多项式凸的紧邻域 L。设 V 是 L 的内部；它是 K 的一个邻域。

依前面的命题，存在一个序列 \((\mathrm{P}_{k})\)：它由 \(C^{n}\) 上以 A 为系数的多项式函数组成，且收敛于 \(f|L\)（依 \(\mathcal{C}(L;A)\) 中的范数）。特别地，序列 \((\mathrm{P}_{k})\) 收敛于 \(f|V\)（在 \(\mathcal{O}(V;A)\) 中）。

按 \(\mathcal{O}(\mathrm{K};\mathrm{A})\) 上拓扑的定义（参见 EVT, II, p. 29, 命题 5），从 \(\mathcal{O}(\mathrm{V};\mathrm{A})\) 到 \(\mathcal{O}(\mathrm{K};\mathrm{A})\) 中的典范映射是连续的。因此，函数 \(\mathrm{P}_k\) 在 K 的某个邻域中的芽所成的序列，收敛于 \(f\) 在 K 的邻域中的芽（在空间 \(\mathcal{O}(\mathrm{K};\mathrm{A})\) 中），这就是所要证明的。

推论 1. --- 设 U 是 K 的一个开邻域。设 \(u_{1}\) 与 \(u_{2}\) 是从 \(\mathcal{O}(\mathrm{U};\mathrm{A})\) 到拓扑空间 X 中的、可经由 \(\mathcal{O}(\mathrm{K};\mathrm{A})\) 分解的连续映射。则 \(u_{1}=u_{2}\) 必须且只需 \(u_{1}\) 与 \(u_{2}\) 在由 \(\mathbf{C}^{n}\) 上以 A 为系数的多项式函数在 K 上的限制所组成的集合上相一致。

推论 2. --- 设 E 是一个 Banach 空间。设 K 是 \(\mathbf{C}^{n}\) 的一个多项式凸紧子集。设 P 是由 \(\mathbf{C}^{n}\) 上取值于 E 的多项式函数在 K 的某个邻域中的芽所组成的集合。则 \(\mathcal{O}(\mathrm{K};\mathrm{E})\) 的每个元素都是 P 中元素所成序列的极限。

给 E 赋以由 ab = 0（对 E 中所有的 a 与 b）定义的乘法 (I, p.~17, 例 1)。这是一个交换 Banach 代数。设 A 是由 E 添加单位元所得到的交换幺 Banach 代数。由于典范映射 \(\mathcal{O}(\mathrm{K};\mathrm{A}) \to \mathcal{O}(\mathrm{K};\mathrm{E})\) 连续，该断言由应用于代数 A 的 Oka-Weil 定理得出。

对 n = 1 与 A = C，我们还有下述结果，它将由 I, p.~150, 推论 2 加以精确化。

定理 3 (Runge). --- 设 K 是 C 的一个紧子集，Q 是在 K 的某个邻域中全纯的有理函数的芽所组成的集合。则 Q 在 \(\mathcal{O}(\mathrm{K})\) 中稠密。

按 \(\mathcal{O}(\mathrm{K})\) 上拓扑的定义，只需证明：对 K 的每个开邻域 U 以及 U 的每个紧子集 L，每个函数 \(f \in \mathcal{O}(\mathrm{U})\) 都是在 L 上连续的有理函数的极限。可以假设 L 是 K 的一个紧邻域。

设 Q' 是 C 上在 L 上连续的有理函数在 L 上的限制所组成的集合，并设 C 是 Q' 在 \(\mathcal{C}(\mathrm{L})\) 中的闭包。这是一个无根的代数。

设 \(z \in \mathrm{C}\) 是 L 的恒等映射。则 C 是 \(\mathcal{C}(\mathrm{L})\) 的由 \(z\) 生成的闭全子代数 (I, p.~6, 引理 2)。因此我们有 \(\mathrm{Sp}_{\mathrm{C}}(z) = \mathrm{Sp}_{\mathcal{C}(\mathrm{L})}(z) = \mathrm{L}\)。于是可以构造 C 中的元素 \(c = \Theta_z(f)\)。由于 C 无根，把 I, p.~66, 推论 1 应用于 C 的特征标 \(g \mapsto g(w)\)（对一切 \(w \in \mathrm{L}\)）就表明 \(c\) 与 \(f\) 在 L 上的限制重合。按 C 的定义，这就证明了 \(f|_{\mathrm{L}}\) 是 \(Q'\) 的元素在 L 上的一个一致极限，从而完成了定理的证明。

